\documentclass[10pt,a4paper]{amsart}
\usepackage[margin=19mm]{geometry}
\usepackage{amsmath,amssymb,mathtools}
\usepackage[scr=rsfs,cal=euler]{mathalfa}
\usepackage{xfrac}
\usepackage[unicode,hidelinks,bookmarks=false]{hyperref}
\newcommand{\ie}{i.e.\ }
\newcommand{\cf}{cf.\ }
\newcommand{\resp}{respectively\ }
\newcommand{\Subsection}{\subsection}
\providecommand{\nobreakdash}{\nobreak\hbox{-}}
\setlength{\emergencystretch}{2em}
\multlinegap0pt
\usepackage{xcolor}
\usepackage{mathtools}
\usepackage{amssymb}
\usepackage{tikz}
\usepackage{float}
\usepackage[shortlabels]{enumitem}
\usepackage{bm}
\usepackage{csquotes}
\newtheorem{lemma}{Lemma}
\usepackage{cleveref}
\crefname{lemma}{Lemma}{Lemmas}
\DeclareMathOperator*{\esssup}{ess\, sup}
\DeclareMathOperator{\inter}{int}
\title{A free boundary approach to the quasistatic evolution of debonding models}
\author{Eleonora Maggiorelli, Filippo Riva, Edoardo Giovanni Tolotti}
\date{Source: arXiv:2503.17023v1 (CC BY 4.0)}
\begin{document}
\maketitle

\noindent\emph{Source and licence.} A free boundary approach to the quasistatic evolution of debonding models. Version arXiv:2503.17023v1 (CC BY 4.0), distributed by arXiv under the Creative Commons Attribution 4.0 International licence, http://creativecommons.org/licenses/by/4.0/, as recorded in the arXiv licence metadata for this exact version and shown on the abstract page https://arxiv.org/abs/2503.17023v1. Full licence notice: \texttt{licenses/arxiv-2503.17023-CC-BY-4.0.txt}.\par\smallskip

\noindent\textit{Editorial scope.} Counted unit: the article's lemma lemma:equivalent (source0.tex lines 408--418). The article's complete original proof of it is delivered with it (source0.tex lines 419--442, one \texttt{proof} environment of 24 lines). No mathematics is written, repaired or re-derived: the statement and the proof are the article's own lines, in the article's own notation, and no retained line is substituted or re-flowed.  Retained uncounted as the source-local context the proof needs: lines 380--386.  Nothing was omitted from any retained span, so the package carries no omission record.  No retained line contains a citation, so the wrapper carries no bibliography and no bibliography entry.  No retained line contains a figure environment, an image-inclusion command or drawing code, so no image is delivered and the package declares no asset dependency.  Editorial formatting only, in addition: an AMS \texttt{amsart} preamble that restates the article's own declarations and macros used by the retained lines, namely the environments \texttt{lemma} on separate counters, together with the standard packages those lines need.  The retained environments print in this document as Lemma 1 in order of appearance, whereas the article prints the same items with the numbering of its own full document.  The complete native source of the article as distributed in the arXiv e-print for this exact version is bundled unchanged as \texttt{sources/175153/source0.tex}.
\begin{lemma}\label{lemma:minimiser_AltCaffarelli}
Let $\eta\in H^1(\Omega)^+$ and let $A\in \mathcal M(\Omega)$. Then the functional $v\mapsto \mathcal{AC}(v,A)$ admits a minimum in the class $H^1_{\Gamma,\eta}(\Omega)$. Moreover, any minimiser $u$ is locally Lipschitz continuous in $\Omega$ and for all $\Omega'$ well contained in $\Omega$ there exists a constant $C'>0$, independent of $A$, such that:
\begin{equation}
    \|u\|_{C^{0,1}(\Omega')}\leq C'(1+\esssup_\Gamma\eta).
\end{equation}
Furthermore, $u$ is harmonic in the open set $\{u>0\}$ and there holds $0\leq u\leq \esssup_\Gamma\eta$ in $\Omega$. Finally, one has $(\inter A)^+_{\Gamma,\eta}\subseteq\{u>0\}$, where $\inter A$ denotes the interior of $A$.
\end{lemma}

\begin{lemma} \label{lemma:equivalent}
    Let $\eta\in H^1(\Omega)^+$ and let $A\in \mathcal M_{\Gamma,\eta}$. For $u \in H^1_{\Gamma,\eta}(\Omega,A)$ the following are equivalent:
    \begin{enumerate}
        \item $\mathcal{AC}(u,A)\le \mathcal{AC}(v,A),\qquad$ for all $v\in H^1_{\Gamma,\eta}(\Omega)$; \label{item:equiv1}
        \item $\mathcal{AC}(u,A)\le \mathcal{AC}(v,A),\qquad$ for all $v\in H^1_{\Gamma,\eta}(\Omega)\cap C^0(\Omega)$;\label{item:equiv2}

        \item $\displaystyle \frac 12\int_\Omega |\nabla u|^2 \, dx  \leq \frac 12\int_\Omega |\nabla v|^2 \, dx + \int_{\{v > 0\}\setminus A}\kappa\, dx,\qquad$ for all $v\in H^1_{\Gamma,\eta}(\Omega)$;\label{item:equiv3}
        \item $\displaystyle \frac 12\int_\Omega |\nabla u|^2 \, dx  \leq \frac 12\int_\Omega |\nabla v|^2 \, dx + \int_{\{v > 0\}\setminus A}\kappa\, dx,\qquad$ for all $v\in H^1_{\Gamma,\eta}(\Omega)^+\cap C^0(\Omega)$ satisfying $A\subseteq \{v>0\}$.\label{item:equiv4}
    \end{enumerate}
    If one of the above holds, then $u=\mathfrak{h}_{A,\eta}$ and furthermore it possesses all the properties stated in Lemma \ref{lemma:minimiser_AltCaffarelli}.
\end{lemma}

\begin{proof} 
    \Cref{lemma:minimiser_AltCaffarelli} yields the equivalence between \eqref{item:equiv1} and \eqref{item:equiv2}. Since $u$ belongs to $H^1_{\Gamma,\eta}(\Omega,A)$ by assumption, the set $\{u>0\}$ is contained in $A$; this implies
    \begin{equation}
       \mathcal{AC}(u,A)=\frac 12 \int_\Omega |\nabla u|^2 \, dx, 
    \end{equation}
    and so \eqref{item:equiv1} is also equivalent to \eqref{item:equiv3}.

    Trivially \eqref{item:equiv3} implies \eqref{item:equiv4}, so we just need to prove the reverse implication; equivalently, we will actually show that \eqref{item:equiv4} implies \eqref{item:equiv2}.

    To this aim, let us fix $v\in  H^1_{\Gamma,\eta}(\Omega)\cap C^0(\Omega)$. By outer regularity of the Lebesgue measure, for any $\varepsilon >0$ there exists an open set $A_\varepsilon\subseteq \Omega$ containing $A$ and satisfying $|A_\varepsilon\setminus A|\le \varepsilon$. We now consider the 1-Lipschitz function $\varphi_\varepsilon(x):=\operatorname{dist}(x;\overline{\Omega}\setminus A_\varepsilon)$, which in particular satisfies $\{\varphi_\varepsilon>0\}=A_\varepsilon$, and we set $v_\varepsilon:=v^++\varepsilon\varphi_\varepsilon $. Observing that $v_\varepsilon\in H^1_{\Gamma,\eta}(\Omega)^+\cap C^0(\Omega)$ and that $A\subseteq A_\varepsilon\subseteq A_\varepsilon\cup \{v>0\}=\{v_\varepsilon>0\}$, by means of \eqref{item:equiv4} we obtain
    \begin{align}
        \mathcal{AC}(u,A)&=\frac 12 \int_\Omega |\nabla u|^2 dx\le \frac 12 \int_\Omega |\nabla v_\varepsilon|^2 dx+\int_{\{v_\varepsilon>0\}\setminus A}\kappa\, dx\\
        &=\frac 12 \int_\Omega |\nabla v^+|^2 dx+\frac{\varepsilon^2}{2}\int_\Omega |\nabla \varphi_\varepsilon|^2 dx+ \varepsilon\int_\Omega \nabla v^+\cdot\nabla \varphi_\varepsilon dx+ \int_{(\{v>0\}\cup A_\varepsilon)\setminus A}\kappa\, dx\\
        &\le \frac 12 \int_\Omega |\nabla v|^2 dx +\int_{\{v>0\}\setminus A}\kappa\, dx+\frac{\varepsilon^2}{2}|\Omega|+\varepsilon|\Omega|^{1/2}\|\nabla v\|_{L^2(\Omega)}+\varepsilon\|\kappa\|_{L^\infty(\Omega)}.
    \end{align}
    By sending $\varepsilon\to 0$ we deduce $\mathcal{AC}(u,A)\le \mathcal{AC}(v,A)$, namely \eqref{item:equiv2} holds true and equivalence between \eqref{item:equiv1}, \eqref{item:equiv2}, \eqref{item:equiv3} and \eqref{item:equiv4} is proved.

    Moreover, we notice that choosing $v=\mathfrak{h}_{A,\eta}$ in \eqref{item:equiv3}, since $\{\mathfrak{h}_{A,\eta}>0\}\subseteq A$ we obtain 
    \begin{equation}
        \frac 12\int_\Omega |\nabla u|^2 \, dx   \leq \frac 12 \int_\Omega |\nabla \mathfrak{h}_{A,\eta}|^2 \, dx \,,
    \end{equation}
    whence $u=\mathfrak{h}_{A,\eta}$  by uniqueness of the minimiser.
    Finally, due to \eqref{item:equiv1}, \Cref{lemma:minimiser_AltCaffarelli} applies to the function $u$.
\end{proof}

\end{document}
